\scanpage{114}
% ms. line 01
\uline{Supposons} maintenant que l'on ait $\sum\alpha_i<+\infty$,
% ms. line 02
posons $f(z)=\prod_1^\infty(1+\alpha_i z)$. Alors $\rho$ est aussi
% ms. line 03
l'ordre de la fonction entière, i.e
% ms. line 04
\[ \rho=\limsup_{0<r\to\infty}\frac{\log(\log f(r))}{\log r} \]

% ms. line 05
En effet, on sait que $\rho$ est $\leq$ ordre de la
% ms. line 06
fonction $f(z)$, \deleted{\unsure{et}} \unsure{montrons|démontrons} \unsure{la|le} contraire. On est
% ms. line 07
ramené à montrer que \deleted{le produit} dans le
% ms. line 08
cas où on a $\alpha_n=\frac1{n^\beta}$, on a
% ms. line 09
$\displaystyle \limsup_{0<r\to\infty}\frac{\log\log f(r)}{\log r}\leq\frac1\beta$ \quad (\unsure{en somme on a|on verra|on sait} en fait $=\frac1\beta$)

% ms. line 10
En effet, on a
% ms. line 11
\[ f(r)=\prod_1^{\unsure{\nu|v}}\left(1+\frac1{n^\beta}r\right)\prod_{\nu+1}^{\infty}\left(1+\frac1{n^\beta}r\right) \]
% ms. line 12
le deuxième facteur est $\leq e^{R(\nu)r}$
% ms. line 13
avec $\displaystyle R(\nu)=\sum_{\nu+1}^\infty\frac1{n^\beta}\leq\int_\nu^\infty\frac{\unsure{dx|du}}{x^\beta}=\frac1{\beta-1}\frac1{\nu^{\beta-1}}$
% ms. line 14
Prenons $\nu$ tel que $r\frac1{\nu^{\beta-1}}\leq r^{\frac1\beta}$, i.e $\nu\geq r^{\frac1\beta}$
% ms. line 15
de façon précise, $\nu=[r^{\frac1\beta}]+1$. Alors, $e^{R(\nu)r}$
% ms. line 16
sera une \uline{fonction de $r$} d'ordre $\leq\frac1\beta$. Il
% ms. line 17
suffit de le vérifier pour le produit $\prod_1^\nu(1+\frac1{n^\beta}r)$
% ms. line 18
$\leq(1+r)^\nu=e^{\nu\log(1+r)}$, qui est manifestement
% ms. line 19
d'ordre $\leq\frac1\beta$.

% ms. line 20
Rappelons que si on pose $f(\unsure{z|r})=\sum a_n z^n$,
% ms. line 21
\deleted{\unsure{dans}} et $B(r)=\mathop{\rm Sup}_n(a_n r^n)$, alors on
% ms. line 22
a aussi \qquad ordre \unsure{à|de} $f=\limsup_{r\to\infty}\frac{\log\log B(r)}{\log r}$
% archive sheet number: 115
